%% kapa: exam=2 round=27 year=2016 period=3 subject=compensation_law subject_ko=감정평가 및 보상법규 seq=1 number=1 points=40
%% source: Q-Net 감정평가사 기출문제 (cst00309) artlSeq=5185401 file=extracted/1008/5185401/제27회 감정평가사 2차시험 기출문제/제27회 감정평가사 2차 시험문제지/(3교시 시험문제지)감정평가 및 보상법규.pdf
%% conversion: kordoc 4.24.51 -> tools/gen2.py (auto); review=proofread-gpt-sol; flags=none
\begin{longquestion}{【문제 1】}
「공익사업을 위한 토지 등의 취득 및 보상에 관한 법률」(이하 ‘토지보상법’이라 함)의 적용을 받는 공익사업으로 인하여 甲은 사업시행자인 한국철도시설공단 乙에게 협의절차를 통해 자신이 거주하고 있던 주거용 건축물을 제공하여 생활의 근거를 상실하게 되었다고 주장하면서 토지보상법 제78조제1항에 따른 이주대책의 수립을 신청하였다. 이에 대해 乙은 “위 공익사업은 선형사업으로서 철도건설에 꼭 필요한 최소한의 토지만 보상하므로 사실상 이주택지공급이 불가능하고 이주대책대상자 중 이주정착지에 이주를 희망하는 자의 가구수가 7호(戶)에 그치는 등 위 공익사업은 토지보상법 시행령 제40조제2항에서 규정하고 있는 이주대책을 수립하여야 하는 사유에 해당되지 아니한다”는 이유를 들어 甲의 신청을 거부하였다. 다음 물음에 답하시오. (40점) \par
1) 乙이 甲에 대한 거부처분을 하기에 앞서 행정절차법상 사전통지와 이유 제시를 하지 아니한 경우 그 거부처분은 위법한가? (20점) \par
2) 만약 甲이 거부처분 취소소송을 제기하였다면, 乙은 그 소송 계속 중에 처분의 적법성을 유지하기 위해 “甲은 주거용 건축물에 계약체결일까지 계속하여 거주하고 있지 아니하였을 뿐만아니라 이주정착지로의 이주를 포기하고 이주정착금을 받은 자에 해당하므로 토지보상법 시행령 제40조 제2항에 따라 이주대책을 수립할 필요가 없다”는 사유를 추가ㆍ변경할 수 있는가? (20점) \par
\begin{framedtext}
참조 조문 \par
<공익사업을 위한 토지 등의 취득 및 보상에 관한 법률> \par
제78조(이주대책의 수립 등) ① 사업시행자는 공익사업의 시행으로 인하여 주거용 건축물을 제공함에 따라 생활의 근거를 상실하게 되는 자(이하 "이주대책대상자"라 한다)를 위하여 대통령령으로 정하는 바에 따라 이주대책을 수립ㆍ실시하거나 이주정착금을 지급하여야 한다. \par
② <이하 생략> \par
<공익사업을 위한 토지 등의 취득 및 보상에 관한 법률 시행령> \par
제40조(이주대책의 수립ㆍ실시) ① <생략> \par
② 이주대책은 국토교통부령으로 정하는 부득이한 사유가 있는 경우를 제외하고는 이주대책대상자 중 이주정착지에 이주를 희망하는 자의 가구 수가 10호(戶) 이상인 경우에 수립ㆍ실시한다. <이하 생략> \par
제41조(이주정착금의 지급) 사업시행자는 법 제78조제1항에 따라 다음 각 호의 어느 하나에 해당하는 경우에는 이주대책대상자에게 국토교통부령으로 정하는 바에 따라 이주정착금을 지급하여야 한다. \par
1. 이주대책을 수립ㆍ실시하지 아니하는 경우 \par
2. 이주대책대상자가 이주정착지가 아닌 다른 지역으로 이주하려는 경우 \par
<행정절차법> \par
제21조(처분의 사전 통지) ① 행정청은 당사자에게 의무를 부과하거나 권익을 제한하는 처분을 하는 경우에는 미리 다음 각 호의 사항을 당사자등에게 통지하여야 한다. \par
1. 처분의 제목 \par
2. 당사자의 성명 또는 명칭과 주소 \par
3. 처분하려는 원인이 되는 사실과 처분의 내용 및 법적 근거 \par
4. 제3호에 대하여 의견을 제출할 수 있다는 뜻과 의견을 제출하지 아니하는 경우의 처리방법 \par
5. 의견제출기관의 명칭과 주소 \par
6. 의견제출기한 \par
7. 그 밖에 필요한 사항 \par
제23조(처분의 이유 제시) ① 행정청은 처분을 할 때에는 다음 각 호의 어느 하나에 해당하는 경우를 제외하고는 당사자에게 그 근거와 이유를 제시하여야 한다. \par
1. 신청 내용을 모두 그대로 인정하는 처분인 경우 \par
2. 단순ㆍ반복적인 처분 또는 경미한 처분으로서 당사자가 그 이유를 명백히 알 수 있는 경우 \par
3. 긴급히 처분을 할 필요가 있는 경우 \par
② 행정청은 제1항제2호 및 제3호의 경우에 처분 후 당사자가 요청하는 경우에는 그 근거와 이유를 제시하여야 한다. \par
\end{framedtext}
\end{longquestion}
